\pdfoutput=1
\documentclass[11pt]{article}

\usepackage{ACL2023}
\usepackage{times}
\usepackage{latexsym}
\usepackage[T1]{fontenc}

\usepackage[utf8]{inputenc}
\usepackage{microtype}
\usepackage{fancyhdr}
\usepackage{natbib}

\pagestyle{fancy}
\fancyhf{}
\fancyhead[LE,RO]{EPFL -- Spring Semester 2026}
\fancyhead[RE,LO]{Modern Natural Language Processing - Open Project Literature Review}
\rfoot{Page \thepage \hspace{1pt}}
\renewcommand{\headrulewidth}{1pt}

%%%%%%%%%%%%%%%%%%% TITLE AND AUTHORS %%%%%%%%%%%%%%%%%%%

\title{Literature Review: Faithfulness and Retrieval Robustness in Retrieval-Augmented Generation}

\author{\normalfont
Elie Bruno | 355932 | \texttt{elie.bruno@epfl.ch} \\
Andrea Trugenberger | SCIPER | \texttt{andrea.trugenberger@epfl.ch} \\
Author 3 | SCIPER 3 | \texttt{author3@epfl.ch} \\
Author 4 | SCIPER 4 | \texttt{author4@epfl.ch} \\
Team Faithful RAG \\
}

%%%%%%%%%%%%%%%%%%% LITERATURE REVIEW %%%%%%%%%%%%%%%%%%%

\begin{document}
\maketitle
\thispagestyle{fancy}

Retrieval-Augmented Generation (RAG) grounds language model outputs in external documents to reduce hallucination without retraining~\cite{retrievalaugmentedgeneration}. Research has since branched in three directions: improving what gets retrieved, verifying whether generation stays faithful to it, and evaluating these systems at scale.

\subsection{Retrieval Methods and Robustness}

Dense passage retrieval~\cite{karpukhin-etal-2020-dense} replaced sparse methods with learned embeddings for semantic matching. However, retrieval quality remains fragile: \citet{shi2023large} showed LLMs degrade when irrelevant context is mixed with relevant passages, and \citet{liu2024lost} demonstrated a positional bias where models neglect information in the middle of long contexts. On the representation side, ColBERTv2~\cite{santhanam2022colbertv2} introduced late-interaction retrieval with token-level similarity, achieving stronger matching than single-vector bi-encoders.

Several approaches tackle retrieval quality directly. Corrective RAG~\cite{yan2024corrective} classifies retrieved documents as correct, ambiguous, or incorrect and triggers query refinement for low-confidence retrievals. Adaptive-RAG~\cite{jeong2024adaptiverag} routes queries by complexity, while FLARE~\cite{jiang2023active} interleaves generation with retrieval by detecting low-confidence spans mid-generation. These reduce noise fed to the generator, but none measures whether this translates to more faithful citations.

\subsection{Faithfulness and Self-Correction}

Self-RAG~\cite{asai2024selfrag} trains the LLM to emit reflection tokens signaling retrieval need, passage relevance, and generation support---internalizing the verify loop but requiring fine-tuning. RAFT~\cite{zhang2024raft} fine-tunes on relevant and distractor documents, teaching models to cite correctly. RankRAG~\cite{yu2024rankrag} unifies reranking and generation in a single instruction-tuned LLM.

For evaluation, FActScore~\cite{min-etal-2023-factscore} decomposes generations into atomic facts verified against a knowledge source, shifting evaluation from surface overlap toward entailment-based verification. ALCE~\cite{gao2023enabling} extends this to citation-specific evaluation, benchmarking whether each generated statement is backed by a correct inline citation.

\subsection{Evaluation Frameworks and Structured Retrieval}

RAGAS~\cite{es2024ragas} defines component-wise metrics: faithfulness, answer relevancy, context precision, and context recall. ARES~\cite{saadfalcon2024ares} automates evaluation via LLM judges with statistical guarantees. The RGB benchmark~\cite{chen2024benchmarking} tests noise robustness, negative rejection, information integration, and counterfactual robustness---axes that overlap with but do not cover citation accuracy. Beyond flat vector search, HippoRAG~\cite{gutierrez2024hipporag} uses knowledge graphs with PageRank traversal, MemoRAG~\cite{qian2024memorag} introduces a dual-system memory architecture, and RECOMP~\cite{xu2024recomp} compresses passages before generation. \citet{gao2024retrievalaugmented} survey these approaches, identifying retrieval granularity, noise filtering, and faithful generation as persistent open problems.

\subsection{Positioning Our Work}

We combine corrective retrieval~\cite{yan2024corrective} with post-hoc citation verification inspired by FActScore and ALCE, without fine-tuning the generator (unlike Self-RAG and RAFT). We move beyond general-domain benchmarks by constructing a domain-specific evaluation set for scientific policy literature. No prior work combines corrective retrieval, claim-level faithfulness scoring, and domain-specific evaluation in a single pipeline.

%%%%%%%%%%%%%%%%%%% BIBLIOGRAPHY %%%%%%%%%%%%%%%%%%%

\newpage
\bibliography{anthology, custom}
\bibliographystyle{acl_natbib}

\end{document}
